\documentclass[11pt]{book}
\usepackage{fontspec}
\setmainfont{TeX Gyre Pagella}
\usepackage[most]{tcolorbox}
\usepackage{listings}
\usepackage{xcolor}
\usepackage{hyperref}

\lstdefinelanguage{Rust}{
  morekeywords={let,mut,fn,struct,enum,impl,trait,match,if,else,for,while,loop,
    return,pub,mod,use,self,Self,crate,super,derive,dyn,move,ref,as,in,where},
  morekeywords=[2]{String,str,Vec,Option,Result,Rc,RefCell,Box,Iterator,i32,u32,i64,f64,bool,char,usize},
  sensitive=true,
  morecomment=[l]{//},
  morestring=[b]",
}
\lstset{
  language=Rust,
  basicstyle=\ttfamily\small,
  keywordstyle=\bfseries,
  keywordstyle=[2]\itshape,
  commentstyle=\itshape,
  showstringspaces=false,
  breaklines=true,
  frame=single,
  framesep=4pt,
  columns=flexible,
}

\newtcolorbox{firstprinciples}{
  colback=gray!5, colframe=black!60, boxrule=0.4pt,
  title={Opening Question}, fonttitle=\bfseries, sharp corners
}
\newtcolorbox{invariantbox}{
  colback=black!3, colframe=black, boxrule=0.6pt,
  title={Invariant Established}, fonttitle=\bfseries, sharp corners
}

\title{Interlude: Iteration Without Materialization}
\author{Thinking Like Rust}
\date{}

\begin{document}

\chapter*{Interlude: Iteration Without Materialization}
\addcontentsline{toc}{chapter}{Interlude: Iteration Without Materialization}

\begin{firstprinciples}
Does a program need to construct a new collection before it can describe what should happen to every element of an existing one?
\end{firstprinciples}

\section*{1. Repetition over a collection}

Begin with the ordinary case.

\begin{lstlisting}
let numbers = vec![1, 2, 3];

for number in &numbers {
    println!("{number}");
}
\end{lstlisting}

Writing \texttt{for number in \&numbers} borrows \texttt{numbers} for the duration of the loop, produces one \texttt{\&i32} per element in turn, and leaves \texttt{numbers} itself fully intact and reachable afterward. Nothing here has been consumed or transformed yet; the loop only observes.

\section*{2. Three ways to enter iteration}

That borrowing choice is one of three, and Rust names all three explicitly rather than leaving the choice implicit.

\begin{lstlisting}
for n in numbers.iter()      { /* &i32   — observe */ }
for n in numbers.iter_mut()  { /* &mut i32 — authorized to mutate */ }
for n in numbers.into_iter() { /* i32    — consumed and transferred */ }
\end{lstlisting}

\texttt{.iter()} grants knowledge of each element without authority over it. \texttt{.iter\_mut()} grants a temporary authorization to mutate each element in place, without transferring ownership of the collection itself. \texttt{.into\_iter()} consumes the collection outright, transferring ownership of each element out of it one at a time; after this loop, \texttt{numbers} no longer exists. These three are not stylistic variants of one operation. They are the same three-way distinction, observation versus mutation-in-place versus consuming transfer, that shows up wherever Rust grants access to a value at all, applied here to what a loop is permitted to do with each element in turn.

\section*{3. The iterator as suspended production}

An iterator is not the collection, and it is not a second collection built to hold the results. It is a rule for producing the next item, together with whatever state is needed to know what ``next'' means the next time it is asked. Calling \texttt{.iter()} on a \texttt{Vec} does not walk the vector; it constructs a small value, an iterator, that has not yet produced anything and will not until something asks it to.

\section*{4. \texttt{next} makes continuation explicit}

That asking has an exact name. Chapter 10 showed a type implementing \texttt{Iterator} by defining \texttt{next}; here, the same method is worth calling directly, from the consuming side, to see what it actually hands back.

\begin{lstlisting}
let mut it = numbers.iter();
println!("{:?}", it.next()); // Some(&1)
println!("{:?}", it.next()); // Some(&2)
println!("{:?}", it.next()); // Some(&3)
println!("{:?}", it.next()); // None
\end{lstlisting}

Each call either advances production and returns \texttt{Some(item)}, or reports that the sequence has ended by returning \texttt{None}. \texttt{Option<Item>} is doing exactly the job Chapter 13 gives it elsewhere: recording, honestly, that ``no further item'' is a distinct and expected outcome rather than a special sentinel value smuggled into the item type itself. Every other iterator method, however elaborate, is ultimately just a disciplined way of calling \texttt{next} repeatedly on the consumer's behalf.

\section*{5. Adapters reshape future production}

An adapter takes an iterator and returns a new iterator describing a changed production rule, without running anything itself.

\begin{lstlisting}
let evens_doubled = numbers
    .iter()
    .filter(|&&n| n % 2 == 0)
    .map(|&n| n * 2);
\end{lstlisting}

\texttt{.filter(...)} describes a rule that skips whatever the inner iterator produces until an item passes the predicate. \texttt{.map(...)} describes a rule that transforms whatever the inner iterator hands it. Neither call has, at this point, looked at a single element of \texttt{numbers}.

\section*{6. Laziness}

That last claim is worth demonstrating rather than asserting, since it's easy to read the chain above and assume the transformation already happened.

\begin{lstlisting}
let chain = numbers.iter().map(|n| {
    println!("touching {n}");
    n * 2
});
println!("chain built, nothing printed yet");
let doubled: Vec<i32> = chain.collect();
// only now does "touching 1" / "touching 2" / "touching 3" appear
\end{lstlisting}

Building \texttt{chain} does not run the closure passed to \texttt{map} even once. The chain is a composed \emph{description} of a possible computation, in exactly the sense Chapter 2's composition argument already established for adapter chains in general; what this example adds is visible proof that the description and its execution are two separate events, separated here by an explicit \texttt{println!} that only fires once something finally drives the iterator forward.

\section*{7. Consumers create consequences}

That driving-forward has its own name too: a consumer. \texttt{sum}, \texttt{count}, \texttt{for\_each}, \texttt{find}, and \texttt{collect} are all consumers — each one repeatedly calls \texttt{next} on whatever chain precedes it until the chain reports \texttt{None}, and produces a final value or effect from what it saw along the way. Nothing before a consumer in a chain has done any work; everything after one is finished.

\section*{8. Why \texttt{collect} needs a destination}

Most consumers have an unambiguous result type: \texttt{sum} produces a number, \texttt{count} produces a \texttt{usize}. \texttt{collect} is different, because many different collection types could equally well receive the same sequence of items.

\begin{lstlisting}
let doubled: Vec<i32> = numbers.iter().map(|n| n * 2).collect();
// or, without a type annotation on the binding:
let doubled = numbers.iter().map(|n| n * 2).collect::<Vec<i32>>();
\end{lstlisting}

The iterator chain itself carries no information about whether its output should become a \texttt{Vec}, a \texttt{HashSet}, or a \texttt{String}; \texttt{collect} is generic over its return type, and something else, a type annotation on the binding or the explicit \texttt{::<Vec<i32>>} turbofish syntax, has to supply the missing piece. This is not a quirk of \texttt{collect} specifically. It is the same parametric-admissibility discipline Chapter 10 formalized: \texttt{collect}'s bound tells the compiler how to build \emph{any} sufficiently well-behaved destination, and the caller still has to say which one is meant.

\section*{9. Short-circuiting and narrowed search}

Because production only happens on demand, a consumer need not exhaust the whole chain to produce its answer.

\begin{lstlisting}
let first_even = numbers.iter().find(|&&n| n % 2 == 0);
\end{lstlisting}

\texttt{find} calls \texttt{next} only until a matching item appears, then stops; if \texttt{numbers} held ten million elements and the first one already matched, the other 9,999,999 would never be produced at all. \texttt{any}, and a bounded \texttt{take} followed by any consumer, work the same way. Laziness is not merely a description of when work happens; here it directly reduces how much work happens, by letting the consumer decide when enough of the sequence has been seen.

\section*{10. Loops and chains are not enemies}

The chain from Section 5 and an ordinary loop describe the same transformation at two levels of compression.

\begin{lstlisting}
// as a chain
let result: Vec<i32> = numbers.iter()
    .filter(|&&n| n % 2 == 0)
    .map(|&n| n * 2)
    .collect();

// as a loop
let mut result = Vec::new();
for &n in numbers.iter() {
    if n % 2 == 0 {
        result.push(n * 2);
    }
}
\end{lstlisting}

Neither form is being ranked above the other here. The chain names each stage as an inspectable value and makes the description-versus-execution split of Section 6 visible in the syntax itself; the loop makes the same transformation explicit as a sequence of statements. Translating between them is a useful exercise precisely because it shows the identical continuation being carried out by two different means, one compressed, one expanded.

\begin{invariantbox}
An iterator separates the description of a traversal from its materialized consequence. Its adapters reshape what may be produced next without producing anything themselves; its consumers decide when that latent traversal becomes actual, and how much of it needs to run at all.
\end{invariantbox}

\bigskip
\noindent\textit{Non-overlap note: Chapter 2 supplies the composition theory behind adapter chains, including the \texttt{Fn}/\texttt{FnMut}/\texttt{FnOnce} treatment of closures, and Chapter 10 supplies the generic-trait theory behind \texttt{Iterator}'s one required method and the functor law for \texttt{map}. This interlude does not re-argue either. It supplies the operational syntax and execution model those two chapters assume but do not separately walk through: \texttt{iter}/\texttt{iter\_mut}/\texttt{into\_iter}, \texttt{next} called from the consuming side, laziness demonstrated rather than asserted, \texttt{collect}'s destination-type problem, and short-circuiting.}

\end{document}
